% Generated from validated synthesis. Do not hand-edit.
\section*{Monthly Signals}
\addcontentsline{toc}{section}{Monthly Signals}
\smallskip
\noindent\textbf{Frontier Models — 性能だけでなく提供形態が競争軸になる} GPT-5.5、Claude Opus 4.7、Qwen、DeepSeek V4、Gemini、GLMの更新は、API lifecycle、互換interface、multimodal／agent接続、open ecosystemへの波及という異なる軸を浮かび上がらせた。 \hfill{\footnotesize p.~\pageref{pkg:frontier-models-april}}\par
\smallskip
\noindent\textbf{Multimodal — 生成・対話・視覚・空間が別々に伸びる} ChatGPT Images 2.0、MiniCPM-o 4.5、Voxtral TTS、GLM-5V-Turboなどを通して、後の統合multimodalへ向かう複数の入口が4月に並行して現れた。 \hfill{\footnotesize p.~\pageref{pkg:multimodal-interfaces-april}}\par
\smallskip
\noindent\textbf{Agents \& Serving — 実行環境化の部品が揃い始める} sandbox、computer use、memory、orchestrationに加え、workflow semantics、KV state、multimodal MoEを意識するserving研究が登場し、modelの外側が設計対象になり始めた。 \hfill{\footnotesize p.~\pageref{pkg:agents-tooling-april} / p.~\pageref{pkg:inference-serving-april}}\par
\smallskip
\noindent\textbf{Evaluation \& Safety — Agentの実行条件を分解して測る} semantic privilege separation、GUI workflow benchmark、PII filtering、visual prompt injection、behavioral signalは、異なるthreat／evaluation modelとして実行境界の設計を前景化させた。 \hfill{\footnotesize p.~\pageref{pkg:evaluation-safety-sandbox-april}}\par
\smallskip
\noindent\textbf{Science \& Paper Watch — 専門化と次の問題設定} GPT-Rosalind、SPUR、DPLM-Evoによる科学領域の専門化と、security・RL rollout・live evaluation・tool-use costを扱う研究が、5月以降へ続く複数の研究課題を示した。 \hfill{\footnotesize p.~\pageref{pkg:ai-science-april} / p.~\pageref{pkg:paper-watch-april}}\par
\medskip
\begin{claimboundary}[Retrospective scope]
本号は2026年4月を後日確認可能になった一次情報も用いて再構成するRetrospective Specialである。Coverage Periodと制作時点を同一視せず、vendor / project / author claimの境界は各記事とTechnical Notesで明示する。
\end{claimboundary}
\medskip
\setcounter{tocdepth}{1}
\tableofcontents
